\documentclass[conference]{IEEEtran}
\usepackage[T1]{fontenc}
\usepackage{cite}
\usepackage{amsmath,amssymb}
\usepackage{graphicx}
\usepackage{booktabs,array}
\usepackage{xcolor}
\usepackage{url}
\usepackage[hidelinks]{hyperref}
\hypersetup{pdftitle={When Operator Accuracy Does Not Transfer to the Image},
pdfauthor={Sally Mohlakola},pdfsubject={Learned interior transport for faceted gemstones}}
\graphicspath{{figures/}}
\newcommand{\dg}{^{\circ}}
\newcommand{\E}{\mathbb{E}}
\newcommand{\R}{\mathbb{R}}
\newcommand{\spp}{\,\mathrm{spp}}
\pagestyle{plain}
\raggedbottom

\begin{document}
\title{When Operator Accuracy Does Not Transfer to the Image:\\
A Learned Interior Transport Operator for Faceted Gemstones}
\author{\IEEEauthorblockN{Sally Mohlakola}
\IEEEauthorblockA{Honours Research Report}}
\maketitle

\begin{abstract}
This study asks whether a compact neural appearance model, as used for
cloth, can reproduce resolved, specular gemstones. A direction-conditioned
radiance distribution map is first evaluated and then generalised to an
interior transport operator that, from entry position, direction, normal and
wavelength, predicts escape, an exit triangle, and an exit position and
direction. Stochastic and deterministic decoders are evaluated on three diamond
cuts against spectral reference renders with noise-corrected measures and a
separation of operator-mediated light from untouched reflections.
Learned exit-triangle prediction exceeds the tested lookup baselines, but this
does not ensure image agreement: stochastic decoders reach pixel correlations of
0.136--0.293, deep deterministic decoders 0.038--0.168 despite higher validation
accuracy. A recorded, convergence-trained run on an independent held-out pool
locates the failure: exit triangles are predicted correctly for 61\% of short
paths but at most 7\% of paths meeting four or more facets, and for only 6\% of
exits through the pavilion. The same run measures cost: each learned interior
decision is 3.5 times faster than tracing it, yet the neural image settles at an
error that a single analytic render already beats. Under the tested studio
illumination, operator-level accuracy and perceived sharpness are insufficient
proxies for rendered fidelity.
\end{abstract}

\begin{IEEEkeywords}
neural rendering, spectral light transport, gemstones, mixture density networks,
image evaluation, Monte Carlo rendering
\end{IEEEkeywords}

\begin{figure*}[t]
\centering
\includegraphics[width=\textwidth]{fig1_overview.png}
\caption{The learned interior operator. Camera rays meet the stone's first
surface, where Fresnel reflection or refraction is traced exactly. A refracted
ray's entry state (position, direction, normal, wavelength; 10 numbers) passes
through a shared two-layer encoder to an escape head, a facet head over the
mesh's triangles, and a 16-number embedding of the chosen facet. One of three
decoders then places the exit on that facet (point and direction), and the exit
ray returns to the path tracer. The deep decoder alone also reads the entry
state through its own trunk.}
\label{fig:overview}
\end{figure*}

\section{Introduction}
The appearance of a faceted diamond depends on where internally reflected and
refracted light emerges. Small changes in entry position or direction can
change the intersected facets and redirect a path. Narrow sources produce
high-contrast structure sensitive to these changes; dispersion adds wavelength
dependence. An approximation must therefore preserve both transmitted energy
and its placement.

Soh and Montazeri~\cite{soh2024} aggregate explicit cloth-fibre scattering into
a compact neural appearance model. This motivates an adaptation to gemstones,
where visible facets and their internal reflections are resolved in the image.
Averaging their angular structure can remove the features the renderer must
reproduce.

The research question is consequently: \emph{how faithfully can a compact
learned interior transport operator reproduce resolved gemstone appearance, and
which operator errors explain its image-space failures?} The study begins with
a radiance distribution map (RDM), diagnoses its missing spatial conditioning,
and evaluates an operator that returns the outcome of an entire interior path
(Fig.~\ref{fig:overview}).

The contributions are fourfold. First, lookup and input-ablation experiments
quantify the benefit of spatial conditioning for exit-triangle prediction.
Second, a factorised operator supports stochastic and deterministic placement
while preserving explicit geometric exit constraints. Third, an evaluation
combining noise-corrected pixel statistics, geometric-facet averages and
path-component separation exposes a mismatch between operator accuracy and
image agreement on three cuts. Fourth, a recorded pipeline (convergence-based
training, an independent held-out pool, per-path-length and per-region
diagnostics, and timing at equal quality and equal time) shows where the
operator fails and what it costs. The conclusions concern the evaluated models,
budgets and scene, rather than all learned representations of specular media.

\section{Background and Related Work}
\textbf{Aggregated appearance.} The cloth model of Soh and
Montazeri~\cite{soh2024} learns an aggregate BSDF from traced fibre-level
transport. The present study adapts the idea of compressing internal
scattering, not the cloth geometry or its complete implementation; its results
investigate this adaptation and are not a benchmark of the original method.

\textbf{Transport through bounded objects.} Vicini et
al.~\cite{vicini2019} learn shape-adaptive subsurface scattering with a
conditional variational autoencoder and local shape descriptors, a precedent for
predicting transport through a region. Here the medium is a smooth dielectric
with discrete reflection/refraction choices and sharply directed exits.
Classical gemstone rendering models the relevant optics~\cite{sun2000}; a
spectral path tracer provides the reference~\cite{pharr2023,wilkie2014}.

\textbf{Distributions, representations and evaluation.} Mixture density
networks~\cite{bishop1994} represent several outcomes for one input, suiting an
entry state that can escape by several internal routes, but do not by themselves
guarantee accurate branch probabilities or placement. Spectral
bias~\cite{rahaman2019} motivates Fourier features~\cite{tancik2020}; hash
encodings~\cite{muller2022ngp} are an alternative. LEAN mapping~\cite{olano2010}
motivates retaining second moments. SSIM and FLIP~\cite{wang2004,andersson2020}
are complementary image criteria; this study instead emphasises noise-corrected
luminance agreement and facet averages, and claims no SSIM or FLIP results.

\section{From a Radiance Map to an Interior Operator}
\label{sec:rdm}
The project began as a direct adaptation of the cloth model to diamond. Its
radiance distribution map (RDM) is gathered by firing rays at the stone,
tracing each one through the interior, and binning the light that leaves by
outgoing direction for each incoming direction. Small networks are fitted to
the binned table, and the result is applied as a material on the stone's
surface. Each step below removed one explanation for why this route did not
produce a diamond, and together they led to the learned interior operator.

\textbf{Is it a sampling problem?} Soh and Montazeri gather 50--100 million
fully traced rays per yarn material and continue until the map
converges~\cite{soh2024}. To give the RDM the same chance, it is gathered for
each cut at five budgets from $2^{26}$ rays (67 million, the scale of the cloth
dataset) to $2^{30}$ (1.07 billion), and each map is rendered under the same
studio rig (Fig.~\ref{fig:ladder}). If the RDM's flat shading were sampling
noise, it would sharpen along this ladder; if it persists while the bins
converge, the limit lies in what the map conditions on rather than in how many
rays fill it. \emph{[Ladder renders pending; this paragraph will report what
changes between $2^{26}$ and $2^{30}$.]}

\begin{figure*}[t]
\centering
\includegraphics[width=0.78\textwidth]{rdm_ladder.pdf}
\caption{RDM renders against gathering budget, after the 50--100 million rays
per material of the cloth model~\cite{soh2024}. For each cut the large panel
uses $2^{30}$ rays and the stacked panels $2^{26}$ to $2^{29}$; all use the same
studio rig and render settings. \emph{Placeholder layout: renders pending.}}
\label{fig:ladder}
\end{figure*}

\textbf{Averaging erases the detail.} Each bin stores the mean of the paths
that land in it. In 99.8\% of well-populated bins the standard deviation
exceeds that mean (median ratio 2.53), so each stored value summarises a
distribution several times wider than itself. More rays make that mean more
precise; they do not make the distribution narrower.

\textbf{Restoring variation is not restoring placement.} A moment-based
modification after LEAN mapping~\cite{olano2010} activates shading points with
probability $p=1/(1+r^2)$ at gain $1/p$, where $r$ is a bin's relative standard
deviation. This preserves the moments in expectation and brings back the
\emph{amount} of sparkle, but the bright points are statistically right and
spatially arbitrary: the result reads as speckle, not as reflections.

\textbf{The missing variable is position.} Asked to predict the exit triangle,
a table conditioned on entry direction, as the RDM is, scores 0.068 with 32
cells and only 0.098 with 2,048, against 0.047 for the most frequent triangle;
conditioning on the entry triangle instead reaches 0.192. Resolution is not the
limit, the conditioning variable is. In yarn, light enters and leaves within a
sub-pixel cross-section, so direction suffices; in a brilliant cut, \emph{where}
light enters decides which facet it leaves through. The learned interior
operator follows from this: it keeps the RDM's idea of replacing the interior
with a learned answer, but conditions on entry position, direction, normal and
wavelength, and returns a single exit rather than a binned average.

\section{Method}
\label{sec:method}
\subsection{Scene, geometry and reference transport}
All comparisons use one studio scene: a perspective camera with a $30\dg$ field
of view, a grey diffuse ground plane, textured area lights and two small intense
sources, i.e.\ \emph{one lighting rig with several sources}. The non-uniform
illumination makes exit-direction errors visible by assigning different radiance
to different parts of the environment.

Rendering uses the scalar spectral variant of Mitsuba 3~\cite{mitsuba3}. The
reference samples Fresnel reflection or refraction, applies wavelength-dependent
refraction from a diamond Sellmeier model, and traces each interior path until
escape or a depth cap of 128, sampling the BSDF without next-event estimation.
Absorption and body colour are not modelled. The \emph{analytic reference} names
these explicit transport rules; a finite-sample image is not asserted to be
converged ground truth.


The three procedural meshes are a round brilliant, pear brilliant and step cut
with 64, 62 and 112 triangles, grouped into 57, 33 and 50 coplanar facets, of
which 25, 17 and 25 are visible in the evaluated view. The round uses crown
angle $34.5\dg$, pavilion angle $40.75\dg$ and table fraction 0.56. Validation
checks closedness, manifold edges, outward orientation and degenerate faces.
The operator classifies \emph{triangles}, which are not optical facets: the
round's table alone is eight triangles. Geometry hashes stored with data and
checkpoints detect mesh mismatches.

\subsection{Conditional exit representation}
The entry vector $\mathbf{x}\in\R^{10}$ comprises position $\mathbf{p}$,
interior direction $\boldsymbol\omega$, entry normal $\mathbf{n}$ and wavelength
$\lambda$, with positions divided by the mesh radius and $\lambda$ encoded as
$(\lambda-595)/235$ (nm). With $e$ escape before the depth cap, $f$ an exit
triangle and $\mathbf{y}\in\R^4$ exit coordinates,
\begin{equation}
 p(e=1,f,\mathbf{y}\mid\mathbf{x})=
 s_\theta(\mathbf{x})\,\pi_\theta(f\mid\mathbf{x})\,
 q_\theta(\mathbf{y}\mid\mathbf{x},f),
\end{equation}
with $p(e=0\mid\mathbf{x})=1-s_\theta(\mathbf{x})$, making non-escape explicit
instead of assigning an artificial exit to a truncated path. Two coordinates
are barycentric logits and two are outgoing tangent slopes:
\begin{equation}
 \boldsymbol\beta=\operatorname{softmax}(y_1,y_2,0),\quad
 \boldsymbol\omega_{\rm out}=
 \frac{\mathbf{n}_f+y_3\mathbf{t}_f+y_4\mathbf{b}_f}
 {\|\mathbf{n}_f+y_3\mathbf{t}_f+y_4\mathbf{b}_f\|},
\end{equation}
with exit point $\mathbf{p}_{\rm out}=\sum_j\beta_j\mathbf{v}_{fj}$. This places
every exit on the selected triangle and in its outward hemisphere, but does not
ensure a physically valid internal route.

\subsection{Model variants and optimisation}
Fig.~\ref{fig:overview} shows the architecture. A two-layer, width-64 SiLU
encoder feeds escape and triangle heads; a learned
16-dimensional triangle embedding conditions the coordinate decoder. The
\emph{stochastic} model uses four diagonal Gaussian mixture
components~\cite{bishop1994}. The \emph{clone} regresses one coordinate vector
with smooth-L1 loss. The \emph{deep clone} uses a four-layer decoder with a
separate encoder. All variants sample the exit triangle; ``deterministic''
refers only to placement given that triangle. On the round cut they have
17,637, 15,557 and 28,741 trainable parameters. Training minimises
\begin{equation}
 \mathcal L_i=\operatorname{BCE}(s_\theta(\mathbf{x}_i),e_i)
 +e_i\left[-\log\pi_\theta(f_i\mid\mathbf{x}_i)
 +\ell_{\rm coord}(\mathbf{y}_i)\right],
\end{equation}
where $\ell_{\rm coord}$ is mixture negative log-likelihood or smooth-L1, so
total losses are not comparable across families. Adam~\cite{kingma2015} uses
learning rate $10^{-3}$, batch 256 and gradient clipping at norm 10.

Two training regimes appear. The historical three-cut runs use fixed budgets
(80 epochs stochastic, 30 deep) and keep the best validation checkpoint. The
\emph{recorded} run (Sec.~\ref{sec:recorded}) instead trains to convergence: it
stops once the best validation objective has improved by no more than $10^{-3}$
for five consecutive epochs, and records training and validation losses, timing,
the command and all input hashes. Decoder comparisons share a training pool per
cut but also change objective and, for the deep clone, capacity; they are model
comparisons, not single-variable ablations.

\subsection{Data and render integration}
Each pool contains 1,048,576 path records from 65,536 entry states drawn to
match the camera's queries over all azimuths. Invalid records are excluded and
all outcomes of an entry state stay on one side of the training/validation
split. For the recorded run a second pool, gathered independently with a
different seed (262,144 records from 16,384 entry states), is used \emph{only}
for testing, so held-out scores are separate from checkpoint selection.

During rendering, entry reflection or transmission follows the analytic Fresnel
rule. Transmitted entries are batched (512) and sent to the operator; each
sampled exit starts a new scene ray, non-escaping events terminate, and exits
that immediately re-enter the stone are discarded. The outside scene and
first-surface reflection use identical code in both modes. The learned
replacement approximates the interior transport and is not an unbiased
estimator of it.

\section{Evaluation Methodology}
\label{sec:evaluation}
\subsection{Sampling, masking and noise baselines}
Renders are $512\times512$ with a box filter. Four independent seeds at
$32\spp$ form a $128\spp$ average per mode; splitting them yields two independent
$64\spp$ halves. Metrics use the mask $M$ of pixels whose primary ray hits the
stone, on linear luminance (Rec.~709 weights); ground caustics lie outside $M$.
The analytic image compared with its own other half gives the \emph{noise
floor}: the agreement sampling noise alone permits at that budget. It is an
empirical baseline, not a strict bound, and does not establish reference
convergence.

\subsection{Image measures}
For neural and reference luminance $N_p,A_p$ the measures are: Pearson
correlation $\rho(N,A)$ of matching halves (does light fall in the right
place?); signed relative energy error
$\epsilon_E=\sum_M N_p/\sum_M A_p-1$ (is the total right?); relative RMSE;
highlight overlap of the brightest 1\% of pixels; and contrast, the standard
deviation over the mean. The recorded run adds four noise-aware measures.
\emph{Disattenuated correlation} $\rho/\sqrt{r_Nr_A}$ divides by each image's
half-split reliability $r$ to estimate agreement without sampling noise.
\emph{Rank} and \emph{log-luminance} correlations check that a score is not
driven by a few bright pixels. \emph{Noise-free contrast} uses the covariance of
independent halves in place of a variance, removing the noise that otherwise
inflates contrast. Ninety-five percent intervals come from bootstrapping
$16\times16$-pixel blocks (pixels) or facets (facet means), which respects
spatial correlation. The measures are read jointly, not combined.

\subsection{Facet averages and component separation}
\label{sec:facets}
Coplanar triangles are grouped into facets; each visible facet covering at
least 20 pixels contributes its mean luminance, and correlation is computed
across facets with equal weight, removing within-facet structure by design.
Each camera path is also tagged by whether it passed through the interior
replacement, so the frame splits exactly as $I=I_{\rm operator}+I_{\rm
untouched}$; \emph{interior-only} scores use $I_{\rm operator}$.

\subsection{Operator diagnostics}
On held-out paths, \emph{top-1 accuracy} asks whether the operator's most
probable triangle is the traced exit triangle; \emph{exit angle} is the error of
the decoded direction when the true triangle is supplied, isolating placement
from selection. The recorded run breaks both down by path length and by exit
region.

\section{Results}
\label{sec:results}
\subsection{Spatial conditioning and learned prediction}
\begin{table}[t]
\caption{Round-cut exit-triangle prediction. Nearest-neighbour results use
4,096 queries over 1,046,359 stored paths.}
\label{tab:baselines}
\centering\small
\begin{tabular}{@{}lr@{}}
\toprule
Predictor & Top-1 accuracy \\
\midrule
Most frequent triangle & 0.047 \\
Entry-direction table (32 / 2,048 cells) & 0.068 / 0.098 \\
Entry-triangle table & 0.192 \\
Nearest neighbour / majority of eight & 0.287 / 0.343 \\
Stochastic operator & 0.380 \\
Clone operator & 0.410 \\
Stochastic, recorded run, independent test pool & 0.374 \\
\bottomrule
\end{tabular}
\end{table}

The learned operators exceed every lookup (Table~\ref{tab:baselines}); the
clone is 43\% above the nearest neighbour. A lookup is a non-parametric
comparator, not an upper bound on learnability. In a 60-epoch conditioning
ablation (seed 23, same split), removing position reduces top-1 from 0.404 to
0.273 and keeping only direction yields 0.117, whereas removing wavelength
changes it by 0.0003. Position therefore matters under this protocol; the small
wavelength effect does not imply that the physical system lacks dispersion.

\subsection{Validation behaviour}
\begin{figure*}[t]
\centering
\includegraphics[width=0.8\textwidth]{validation_curves.pdf}
\caption{Validation histories of the historical runs. Solid: stochastic decoder;
dashed: deep deterministic decoder. Triangle accuracy and cross-entropy are
comparable across families; total losses are not. Single training seed.}
\label{fig:validation}
\end{figure*}



Learning improves triangle prediction on all cuts (Fig.~\ref{fig:validation}),
and deep models reach higher validation accuracy than their stochastic
counterparts: 0.392 against 0.371 on the round, 0.338 against 0.311 on the pear
and 0.328 against 0.305 on the step cut, with lower triangle cross-entropy. Several selected
checkpoints fall on the final epoch, so these fixed-budget curves do not
establish convergence. The recorded run removes this doubt for the round
stochastic model (Fig.~\ref{fig:training}): trained to the convergence
criterion it stops at epoch 62 (best 57, 777\,s) with validation accuracy
0.375, within 0.004 of the fixed-budget model. Validation loss flattens while
training loss keeps falling, so further epochs would mainly fit the training
pool; the network is not under-trained in this configuration.

\begin{figure}[t]
\centering
\includegraphics[width=\columnwidth]{training_curves.pdf}
\caption{Recorded run (round, stochastic): training and validation loss and
validation exit-facet accuracy per epoch. Dashed: the best epoch; training
stopped five epochs later when validation loss had improved by no more than
$10^{-3}$.}
\label{fig:training}
\end{figure}

 The round stochastic model's mean predicted escape probability (0.9984)
also matches the measured escape fraction (0.9975), although with rare
non-escapes this is agreement of rates rather than of individual events.

\subsection{Sharper images have lower pixel agreement}
\begin{figure*}[t]
\centering
\includegraphics[width=\textwidth]{gia_three_heads.png}
\caption{Round-cut renders: stochastic, clone, deep clone and analytic
reference at 128 spp. Deterministic placement restores sharp features but not
their locations. Embedded scores use all seeds; Table~\ref{tab:images} uses
64-spp half-splits.}
\label{fig:round}
\end{figure*}

\begin{table*}[t]
\caption{Matched three-cut image evaluation (historical runs). Four 32-spp seeds
per mode; pixel correlation, relative RMSE, highlight overlap and energy use
64-spp halves; contrast and facet correlation (across visible facet means, each
weighted equally) use the full average. ``Reference halves'' is the noise
baseline.}
\label{tab:images}
\centering\small
\begin{tabular}{@{}llrrrrrr@{}}
\toprule
Cut & Decoder/comparison & Pixel corr. $\uparrow$ & Rel. RMSE $\downarrow$ & Top-1\% overlap $\uparrow$ & Energy error & Contrast & Facet corr. $\uparrow$ \\
\midrule
Round & Stochastic & 0.152 & 3.750 & 0.197 & $-0.8\%$ & 0.93 & 0.399 \\
 & Clone & 0.048 & 4.741 & 0.071 & $+26.4\%$ & 2.35 & 0.288 \\
 & Deep clone & 0.038 & 4.849 & 0.117 & $+13.8\%$ & 2.74 & 0.038 \\
 & Reference halves & 0.957 & 1.140 & 0.777 & $+0.6\%$ & 3.78 & 0.9996 \\
\midrule
Pear & Stochastic & 0.136 & 1.834 & 0.007 & $+17.3\%$ & 0.89 & 0.645 \\
 & Deep clone & 0.054 & 3.795 & 0.002 & $+47.8\%$ & 2.31 & 0.414 \\
 & Reference halves & 0.899 & 0.684 & 0.595 & $+0.4\%$ & 1.48 & 0.9996 \\
\midrule
Step & Stochastic & 0.293 & 2.353 & 0.059 & $+14.2\%$ & 0.84 & 0.880 \\
 & Deep clone & 0.168 & 2.956 & 0.057 & $+15.3\%$ & 1.77 & 0.899 \\
 & Reference halves & 0.942 & 0.817 & 0.650 & $-0.3\%$ & 2.37 & 0.9997 \\
\bottomrule
\end{tabular}
\end{table*}

\begin{figure*}[t]
\centering
\includegraphics[width=\textwidth]{operator_breakdown.pdf}
\caption{Recorded run, 261,632 escaped held-out paths from the independent
test pool. (a) Exit-triangle top-1 accuracy by the number of facets a path meets
inside the stone; the second tick line is each group's share of paths (the
1-facet group has only 82 paths, grey). (b) Angle between the decoded and traced
exit when the true triangle is supplied: median (dot) and 90th percentile.
(c) Share of exits through each region that the operator also sends there; bar:
whole region, dot: one triangle (8 table, 24 crown, 24 of 32 pavilion triangles
with at least 20 test exits; horizontal spread is for visibility only). The
second tick line is the region's share of all exits.}
\label{fig:breakdown}
\end{figure*}

Under the tested rig, deterministic decoders produce sharper, higher-contrast
images but lower pixel correlation and higher relative RMSE on every cut
(Fig.~\ref{fig:round}, Table~\ref{tab:images}): round 0.038 against 0.152, pear
0.054 against 0.136, step 0.168 against 0.293, while reference halves agree at
0.899--0.957. The stochastic round image nearly preserves energy, yet its
contrast is 0.93 against 3.78: energy agreement does not imply structure.
Conversely, the pear deep clone exceeds reference contrast (2.31 against 1.48)
with a $+47.8\%$ energy error, so a sharper image is not necessarily better in
placement or radiometry. Highlight overlap stays far below the reference's own,
independently confirming misplaced bright structure.

\begin{figure}[t]
\centering
\includegraphics[width=\columnwidth]{pear_heads.png}\\[2pt]
\includegraphics[width=\columnwidth]{step_heads.png}
\caption{Pear (top) and step (bottom): stochastic and deep-clone results beside
their references (128 spp). Broad shape and brightness organisation survive
even when fine internal reflections are incorrect.}
\label{fig:cuts}
\end{figure}

\subsection{Coarse agreement is a distinct result}

Facet averaging changes the interpretation (last column of Table~\ref{tab:images}). Stochastic
facet correlations are 0.399, 0.645 and 0.880 for round, pear and step; the deep
clone's are 0.038, 0.414 and 0.899, against reference self-agreement above
0.9995. The step deep clone thus reproduces coarse facet brightness well despite
low pixel correlation. Operator improvement does not worsen \emph{every}
statistic: the effect depends on spatial scale, and a facet mean intentionally
discards the sparkle whose placement fails.

\subsection{Branch selection and within-branch placement}
Within a recorded exit triangle, stochastic samples spread by a median of about
$26.0\dg$, whereas traced branches are almost collimated. A clone has no spread,
but its mean direction error given the true triangle is still about $24.7\dg$:
determinism removes blur without establishing placement. Decoding the same
stochastic checkpoint at component means cuts spread from $26.0\dg$ to
$8.1\dg$, yet image correlation falls from 0.174 to 0.153.

\begin{figure}[t]
\centering
\includegraphics[width=0.88\columnwidth]{oracle_facet.png}
\caption{Clone, clone supplied with the traced exit triangle, and reference
(256 spp). The correct triangle restores contrast but not placement.}
\label{fig:oracle}
\end{figure}

Supplying the clone with the traced exit triangle restores contrast (3.73
against 3.74) but correlation stays at 0.041 (Fig.~\ref{fig:oracle}): the
diamond \emph{looks} right but its highlights land in the wrong places, which
implicates within-triangle placement even when selection is bypassed.

\subsection{A recorded, converged run}
\label{sec:recorded}
To test the evaluation pipeline end to end, one configuration (round cut,
stochastic decoder, seed 23) was re-run from gathered data to final figures with
every input and setting recorded. It is a single configuration: the same
pipeline is set up for all three cuts, both decoders and three training seeds,
and those runs are pending.

\begin{table}[t]
\caption{Recorded run: image agreement with the analytic reference (round cut,
stochastic decoder, $4\times32\spp$). Brackets: 95\% block-bootstrap interval;
``floor'': the analytic image against its own other half.}
\label{tab:recorded}
\centering\small
\setlength{\tabcolsep}{2.5pt}
\begin{tabular}{@{}lrr@{}}
\toprule
Measure & Whole stone & Interior only \\
\midrule
Pixel correlation & 0.141 [0.090, 0.235] & 0.077 [0.063, 0.107] \\
\quad noise floor & 0.957 & 0.959 \\
Disattenuated correlation & 0.213 & 0.156 \\
Rank / log correlation & 0.32 / 0.33 & 0.21 / 0.23 \\
Facet correlation & 0.352 [0.03, 0.90] & 0.246 [0.03, 0.76] \\
Energy error & $-1.7\%$ & $-1.9\%$ \\
Noise-free contrast (ref.) & 0.77 (3.74) & 0.59 (4.32) \\
\bottomrule
\end{tabular}
\end{table}

\begin{figure}[t]
\centering
\includegraphics[width=0.85\columnwidth]{quality_vs_time.pdf}
\caption{Error against render time for the recorded run (whole stone). Solid:
rendered seeds; dashed: the error model's extrapolation (analytic: variance
falls as $1/k$; neural: a fixed bias plus variance$/k$).}
\label{fig:quality}
\end{figure}

\textbf{The earlier result reproduces.} Table~\ref{tab:recorded} matches the
historical round stochastic row of Table~\ref{tab:images} (correlation 0.141
against 0.152, inside the new interval). Removing sampling noise
(disattenuated 0.213) or the influence of bright pixels (rank 0.32, log 0.33)
raises the score but leaves it far below the floor, so the disagreement is
structural rather than noise or a few highlights. The interior-only correlation
(0.077) is roughly half the whole-stone value: the light that actually passes
through the network agrees least, and untouched reflections lift the
whole-stone score. The facet interval is wide because only 25 facets are
visible, so facet scores from a single view should be read with caution.


\textbf{Where the operator fails.} On the independent test pool the operator
chooses the correct exit triangle for 37.4\% of paths, matching its validation
score (37.5\%), so it generalises to new entry states rather than memorising the
training pool. Fig.~\ref{fig:breakdown} shows that this average hides two
regimes. Paths meeting two or three facets (58\% of all paths) are predicted
correctly 61\% of the time and, given the right triangle, leave within a median
$8\dg$ of the traced direction. For paths meeting four or more facets (42\%),
accuracy collapses to 2--7\% and the median exit error grows to
$28\dg$--$46\dg$. By region, exits through the table are predicted correctly
58\% of the time and through the crown 38\%, but through the pavilion only 6\%,
although the pavilion carries a quarter of all exiting light. The crown splits
into two groups of triangles near 55\% and 15--29\%, the lower group having
half the training exits. The pavilion's failure, however, is not a data
shortage: its large triangles each have about $1.6\times10^4$ training exits yet
score 2--13\%. The operator is therefore reliable for short routes out of the
top of the stone and fails for long, many-bounce routes, which are exactly the
paths whose exits depend most sensitively on entry state.

\subsection{Cost at equal time and equal quality}
\begin{table}[t]
\caption{Recorded cost, round cut, $512\times512$, 32 spp per render, mean of
four seeds (four renders in parallel on a 16-thread AMD CPU). Error is relative
RMSE against the converged analytic image (Sec.~\ref{sec:recorded}).}
\label{tab:cost}
\centering\small
\setlength{\tabcolsep}{4pt}
\begin{tabular}{@{}lrr@{}}
\toprule
& Analytic & Neural \\
\midrule
Render time (s) & \textbf{538} & 737 \\
Time per interior decision ($\mu$s) & 63 & \textbf{18} \\
Interior share of render time & 49\% & 10\% \\
Error after four renders & \textbf{0.56} & 3.69 \\
Error at equal time (2,950 s) & \textbf{0.48} & 3.69 \\
Lowest reachable error & $\to 0$ & 3.65 \\
Time to reach the other's best & \textbf{51 s} & never \\
Training time & -- & 777 s \\
Model size & 1 KB mesh & 70 KB \\
\bottomrule
\end{tabular}
\end{table}


\begin{figure}[t]
\centering
\includegraphics[width=0.84\columnwidth]{timing.pdf}
\caption{Where the time goes (recorded run). (a) Render time split into
interior decisions (dark) and everything else (light), with the seed range.
(b) Cost of one interior decision against batch size: a standalone benchmark
(curve, 10th--90th percentile band), the recorded run's benchmark (squares,
different machine load), the analytic trace in the renderer (dashed) and the
network in the renderer at its mean batch (diamond).}
\label{fig:timing}
\end{figure}

Table~\ref{tab:cost} and Figs.~\ref{fig:timing} and~\ref{fig:quality} answer
whether the operator saves time. Error is measured against the converged analytic image: the
analytic error is its sampling variance, which falls with every extra render,
while the neural error is a fixed bias (from the mismatch in
Table~\ref{tab:recorded}) plus a smaller variance. Each learned interior
decision is 3.5 times faster than tracing it, and the interior's share of the
render drops from 49\% to 10\% (Fig.~\ref{fig:timing}a). Yet a full neural
render is 37\% \emph{slower}: the time outside the interior grows from 272\,s to
661\,s, spent gathering entry states into batches, preparing their inputs and
decoding the returned exits. Batch size explains why the per-decision saving is
modest (Fig.~\ref{fig:timing}b). Queried one at a time the network costs about
0.7\,ms, ten times more than tracing; it only overtakes tracing beyond about ten
queries per batch, and reaches 3.5--5\,$\mu$s at 512. The renderer asks for
batches of 512, but as paths terminate its batches average only 80 queries, so
each decision costs 18\,$\mu$s rather than 4. More importantly, the neural error levels off at 3.65, whereas a single analytic
render already reaches 1.12 and the error model places the analytic image at
the neural floor after about 51\,s. At equal time the analytic error is 7.7 times
lower. There is consequently no equal-quality time for the neural method: it
cannot reach analytic quality at any budget. Historical runs agree, with neural
renders 1.2--2.0 times slower than analytic ones across the three cuts.

\subsection{Entry sampling and coverage}
\label{sec:coverage}
\textbf{The aim offset began with the RDM.} The first RDM gathers fired every
ray at the stone's centre. Aiming at one point caps the incidence angle each
facet can receive, so downward-tilted facets never saw some incoming
directions: no checkpoint up to run~12 filled more than 41 of its 64
incident-direction bins, and re-gathering with a centre aim reproduced 41/64.
Firing each ray instead from a random point on a disc the size of the stone's
silhouette, the \emph{aim offset}, together with stratified directions (uniform
in $\cos\theta$ within each angular bin, so grazing directions are
represented), fills all 64 bins in every later checkpoint. The current RDM
gather keeps both. The lesson, that \emph{where} a
ray is aimed matters as much as its direction, carried over to the operator.

\begin{figure}[t]
\centering
\includegraphics[width=\columnwidth]{sampling_contrast.pdf}
\caption{Entry sampling on the pear. Top: where training entries land under
(a) uniform directions with the aim offset and (b) camera rays from a full
orbit. Bottom: a pavilion-facing tumble pose (frame 53 of 60) rendered by
operators trained on (b) alone and on (a) and (b) together.}
\label{fig:sampling}
\end{figure}

\textbf{Camera sampling misses the pavilion.} The operator's gather offers two
strategies (Fig.~\ref{fig:sampling}). Camera sampling matches what the renderer
will ask, but a camera above the stone sees only the crown, so no entries land
on the pavilion; uniform sampling with the RDM's aim offset puts 54\% there. In
a 60-frame pear tumble, an operator trained on camera entries alone lets only
46\% of queries escape when the pavilion faces the camera, and the pavilion
renders nearly black (Fig.~\ref{fig:sampling}c). Adding uniform aim-offset entries to the camera
pool restores escape to 100\% at every pose and lights the pavilion
(Fig.~\ref{fig:sampling}d). Placement, however, is not repaired: correlation at
that pose stays at 0.018 against a floor of 0.951. Coverage is necessary but not
sufficient, and the recorded run, trained on camera entries only, should be
read as a front-view result.

\textbf{Transfer.} A pear-trained operator retargeted to an unseen round
variant by matching facet normals scored 0.37, yet 90\% of its dominant
highlight was untouched light: removing that patch drops correlation to
$-0.001$, and the operator component scores 0.004. Shared reflections can
manufacture agreement, hence the interior-only scores of
Table~\ref{tab:recorded}.

\section{Discussion}
\subsection{Why predictive gains need not improve the image}
For entry state $\mathbf{x}$ and exit state $z$, the image depends on
\begin{equation}
 L(\mathbf{x})=\int L_{\rm ext}(z)\,p(z\mid\mathbf{x})\,\mathrm{d}z.
\end{equation}
An objective that weights all exits alike need not prioritise those that reach
narrow bright sources, where a small angular error matters most; triangle
accuracy and mean angular error are thus incomplete proxies. A broad predicted
distribution spreads energy over correct and incorrect regions, keeping some
correlation while suppressing contrast; a sharper but misplaced one loses that
overlap. This is consistent with the results but not a universal law, as the
step-cut facet result shows.

The recorded breakdown sharpens the account: accuracy depends less on how much
data a triangle has than on how long the route to it is. Many-bounce routes
change triangle whenever a path crosses any internal facet edge, so their exit
varies rapidly with entry state, and smooth densities blend across such
changes~\cite{rahaman2019}. Local encodings~\cite{muller2022ngp} and
radiance-weighted objectives are hypotheses for further work. Incomplete optimisation is a weaker alternative explanation for the
round stochastic model, which converges to the fixed-budget result. The study is
thus an empirical failure analysis of one compact learned replacement, not a
proof that neural transport cannot represent specular gemstones.

\section{Limitations and Further Validation}
\label{sec:limits}
The study uses one lighting rig, mostly one pose, simplified meshes and
mesh-specific triangle heads. The recorded pipeline has run for one
configuration; the historical results use single training seeds and fixed
budgets. Multi-seed runs for all cuts and decoders, and the RDM ladder
(Fig.~\ref{fig:ladder}), are pending. Renders ran four at a time, so times
include contention; reference convergence beyond $128\spp$ is modelled, and
ground caustics are not evaluated.

\section{Conclusion}
A direction-conditioned radiance map was generalised to a position-conditioned
interior transport operator. It beats the tested lookups and generalises to an
independent test pool, but higher accuracy does not bring better fine-scale
image agreement. A recorded run locates the failure in long internal routes and
pavilion exits, and shows the operator faster per decision yet slower per image,
with a bias no sample budget removes. Stronger claims await the multi-seed runs.

\begin{thebibliography}{99}
\bibitem{soh2024} G. Y. Soh and Z. Montazeri, ``Neural appearance model for cloth
rendering,'' \emph{Computer Graphics Forum}, vol. 43, no. 4, 2024,
doi: 10.1111/cgf.15156.
\bibitem{vicini2019} D. Vicini, V. Koltun, and W. Jakob, ``A learned shape-adaptive
subsurface scattering model,'' \emph{ACM Trans. Graph.}, vol. 38, no. 4,
Art. 126, 2019, doi: 10.1145/3306346.3322974.
\bibitem{sun2000} Y. Sun, F. D. Fracchia, and M. S. Drew, ``Rendering diamonds,''
in \emph{Proc. 11th Western Computer Graphics Symposium}, 2000.
\bibitem{pharr2023} M. Pharr, W. Jakob, and G. Humphreys,
\emph{Physically Based Rendering: From Theory to Implementation}, 4th ed.
MIT Press, 2023. [Online]. Available: \url{https://pbr-book.org/4ed/}
\bibitem{wilkie2014} A. Wilkie, S. Nawaz, M. Droske, A. Weidlich, and J. Hanika,
``Hero wavelength spectral sampling,'' \emph{Computer Graphics Forum}, vol. 33,
no. 4, pp. 123--131, 2014, doi: 10.1111/cgf.12419.
\bibitem{bishop1994} C. M. Bishop, ``Mixture density networks,'' Aston University,
Tech. Rep. NCRG/94/004, 1994.
\bibitem{rahaman2019} N. Rahaman \emph{et al.}, ``On the spectral bias of neural
networks,'' in \emph{Proc. 36th Int. Conf. Machine Learning}, PMLR, vol. 97,
pp. 5301--5310, 2019.
\bibitem{tancik2020} M. Tancik \emph{et al.}, ``Fourier features let networks learn
high frequency functions in low dimensional domains,'' in \emph{Advances in
Neural Information Processing Systems}, vol. 33, pp. 7537--7547, 2020.
\bibitem{muller2022ngp} T. M\"uller, A. Evans, C. Schied, and A. Keller,
``Instant neural graphics primitives with a multiresolution hash encoding,''
\emph{ACM Trans. Graph.}, vol. 41, no. 4, Art. 102, 2022,
doi: 10.1145/3528223.3530127.
\bibitem{olano2010} M. Olano and D. Baker, ``LEAN mapping,'' in \emph{Proc. ACM
SIGGRAPH Symp. Interactive 3D Graphics and Games}, pp. 181--188, 2010.
\bibitem{wang2004} Z. Wang, A. C. Bovik, H. R. Sheikh, and E. P. Simoncelli,
``Image quality assessment: From error visibility to structural similarity,''
\emph{IEEE Trans. Image Process.}, vol. 13, no. 4, pp. 600--612, 2004,
doi: 10.1109/TIP.2003.819861.
\bibitem{andersson2020} P. Andersson, J. Nilsson, T. Akenine-M\"oller,
M. Oskarsson, K. \AA str\"om, and M. D. Fairchild,
``FLIP: A difference evaluator for alternating images,''
\emph{Proc. ACM Comput. Graph. Interact. Tech.}, vol. 3, no. 2,
Art. 15, 2020, doi: 10.1145/3406183.
\bibitem{mitsuba3} W. Jakob \emph{et al.}, ``Mitsuba 3 renderer,'' 2022.
[Online]. Available: \url{https://mitsuba-renderer.org}
\bibitem{kingma2015} D. P. Kingma and J. Ba, ``Adam: A method for stochastic
optimization,'' in \emph{Proc. Int. Conf. Learning Representations}, 2015.
\end{thebibliography}
\end{document}
